\capituloSinNumero{Métodos de investigación}
\label{cap:metodos}

Se aplican los métodos teóricos y los métodos empíricos propios de este tipo
de estudio. El análisis y la síntesis relacionan las necesidades del negocio
con los requisitos del sistema: mediante ellos se descompone el proceso de
atención en sus actividades y reglas, y se recomponen en los requisitos de
cupos, reservas, pagos y avisos que la propuesta debe cumplir. La revisión
bibliográfica fundamenta la gestión de reservas turísticas, los métodos de
trabajo y el diseño web, y aporta la comparación con las plataformas
comerciales y con los trabajos académicos citados en los antecedentes.

Los métodos empíricos son la observación, las entrevistas y la revisión
autorizada de documentos. La observación de la atención diaria y las
entrevistas con el personal permiten describir cómo se consultan los cupos y
cómo se confirman las reservas; la revisión autorizada de documentos
internos aporta los datos de la población y las prácticas de registro. La
aplicación de estos métodos se documenta durante el diagnóstico: cada
dificultad identificada queda asociada a una regla de negocio o a un criterio
de aceptación de la propuesta.

La investigación tiene un enfoque mixto. Es cualitativa porque describe los
procedimientos del proceso de venta, las reglas de confirmación y la
percepción del personal mediante entrevistas y observación. Es cuantitativa
porque cuantifica indicadores del alcance y del trabajo: las 58 personas de
la población, los 55 puntos de historia del backlog y el ritmo de avance
medido en el tablero.

El tipo de investigación es descriptivo-propositivo. Descriptivo, porque
caracteriza el proceso actual de venta de paquetes turísticos y sus
dificultades; propositivo, porque propone un sistema web que mejora la
consistencia del proceso de reservas. El diseño corresponde a una
investigación de campo, porque la información se recoge directamente de los
trabajadores y de los clientes de la agencia, y no de fuentes secundarias.
